\chapter{Graphics}
\label{ch:graphics}

\standfirst{One bit per pixel, one operation, and a display that is an
ordinary object you can draw on from a workspace.}

\section{The pieces}

\begin{center}
\small
\begin{tabular}{@{}lp{0.60\textwidth}@{}}
\toprule
\ctp{Point} & \ctp{x @ y}. Arithmetic works: \ctp{(3@4) + (1@1)}\\
\ctp{Rectangle} & \ctp{origin corner:}, \ctp{origin extent:}\\
\ctp{Form} & a bitmap: \ctp{bits width height offset}\\
\ctp{DisplayScreen} & the screen, a Form. The global is \ctp{Display}\\
\ctp{BitBlt} & the one operation that moves rectangles of bits\\
\ctp{Pen} & turtle graphics over a Form\\
\ctp{Cursor} & a 16$\times$16 Form with a hot spot\\
\ctp{StrikeFont} & glyphs, as one wide Form plus an \ctp{xTable}\\
\bottomrule
\end{tabular}
\end{center}

\section{Points and rectangles}

\begin{st}
3 @ 4                              "a Point"
(3 @ 4) + (10 @ 10)                "13@14"
(3 @ 4) x                          "3"

100 @ 100 corner: 300 @ 200        "a Rectangle by two corners"
100 @ 100 extent: 200 @ 100        "the same one, by size"
aRectangle width
aRectangle center
aRectangle containsPoint: 150 @ 150
(aRectangle intersect: anotherRectangle)
\end{st}

\section{Forms}

A \ct{Form} is a width, a height, an offset, and a word array of bits. One bit
per pixel: 1 is black, 0 is white.

\begin{st}
| f |
f := Form extent: 64 @ 64.
f fillBlack.
f displayAt: 100 @ 100.            "draw it on the Display"
\end{st}

\ct{Form} inherits from \ct{DisplayMedium}, which is where the drawing
convenience lives:

\begin{st}
aForm fillBlack.
aForm fillWhite.
aForm fill: aRectangle mask: Form gray.
aForm reverse.                     "invert everything"
aForm border: aRectangle width: 2.
\end{st}

The standard halftone masks are class-side on \ct{Form}: \ct{Form black},
\ct{Form white}, \ct{Form gray}, \ct{Form lightGray}, \ct{Form darkGray},
\ct{Form veryLightGray}.

\section{BitBlt: the one operation}

Everything drawn in Smalltalk-80---every character, every window, every menu,
the rubber band you drag out a window with---is one operation: copy a
rectangle of bits from a source Form to a destination Form, through a halftone
mask, combining with a rule.

\begin{st}
(BitBlt
    destForm: Display
    sourceForm: aForm
    halftoneForm: nil
    combinationRule: Form over
    destOrigin: 100 @ 100
    sourceOrigin: 0 @ 0
    extent: aForm extent
    clipRect: Display boundingBox) copyBits
\end{st}

The combination rules are the sixteen boolean functions of two bits. The ones
with names:

\begin{center}
\small
\begin{tabular}{@{}lll@{}}
\toprule
name & number & effect\\
\midrule
\ctp{Form over} & 3 & source replaces destination\\
\ctp{Form under} & 7 & source ORed in --- black wins\\
\ctp{Form reverse} & 6 & XOR --- doing it twice undoes it\\
\ctp{Form erase} & 4 & clear where the source is black\\
\ctp{Form and} & 1 & AND\\
\bottomrule
\end{tabular}
\end{center}

\ct{Form reverse} is worth remembering because of what it makes possible:
drawing the same rectangle twice with rule 6 leaves the screen exactly as it
was. That is how 1983 draws a rubber band---as a flash between two frames,
never as a mark on the Form.

\section{Pen}

Turtle graphics, and the friendliest way to put something on the screen:

\begin{st}
| p |
p := Pen new.
p defaultNib: 2.
p home.
4 timesRepeat: [p go: 100; turn: 90].
\end{st}

\begin{center}
\small
\begin{tabular}{@{}ll@{}}
\toprule
\ctp{home} \ctp{place:} \ctp{goto:} & position\\
\ctp{go:} \ctp{turn:} \ctp{north} \ctp{direction:} & move\\
\ctp{up} \ctp{down} & lift and lower the nib\\
\ctp{defaultNib:} \ctp{black} \ctp{white} & the nib\\
\bottomrule
\end{tabular}
\end{center}

\section{The Display}

\ct{Display} is a \ct{DisplayScreen}, which is a \ct{Form}, which means the
screen is an ordinary object and you can draw on it from a workspace:

\begin{st}
Display fillWhite.
Display fillBlack: (100 @ 100 extent: 50 @ 50).
Display border: (10 @ 10 corner: 200 @ 100) width: 3.
Display extent.
ScheduledControllers restore.      "put the windows back"
\end{st}

That last line is how you clean up: the window system redraws everything it
knows about.

\begin{stnote}[The screen is the window]
The image asks \ct{Display extent} every time it wants to know how big the
screen is, and nothing caches it---so rather than letterbox a fixed 640$\times$480
screen into your window, this system \emph{grows the Form to fill the window}.
The Form's identity is kept, because \ct{Display} is reachable from every
controller and view; only the three fields change. That is why the desktop
fills a tiling window manager's tile with no border.
\end{stnote}

\section{How antialiased text is possible}

The display is one bit per pixel and has to stay that way: BitBlt is checked
byte-for-byte against the Xerox traces, and every rule is defined on single
bits. Giving \ct{Form} a depth would be a different system.

So the smooth text is not in the image at all. There is a second plane---ink
coverage from 0 to 255---maintained by \emph{recognising text}: a blit whose
source is exactly the font strike is a character being drawn, and the source
$x$ says which character, because that is precisely what the \ct{xTable}
means. The same glyph is stamped from an eight-bit coverage table at the same
advance, so the two agree by construction.

Anything else that lands on the display \emph{drops} the coverage under it.
Scroll a pane, clear a window, invert a selection, and the shadow gives up and
the one-bit pixels show through; text comes back smooth the moment it is drawn
again. That conservatism is what keeps it honest---modelling what every rule
does to coverage is how you end up with a second graphics system that
disagrees with the first.

Nothing in the image can observe any of this. \ct{Display} is one bit deep to
every Smalltalk method that asks.

\section{Views, if you want to build one}

The windowed interface is Model--View--Controller, and it is the 1983 version:
a \ct{Model} holds state and broadcasts \ct{changed:}; a \ct{View} draws it; a
\ct{Controller} handles input. Views nest, each with its own transformation
and clipping rectangle.

Building a new tool means subclassing \ct{View} and \ct{Controller}, or more
usually \ct{StringHolderView} and \ct{StringHolderController}, which already
do text. \ct{StandardSystemView} is what makes something a window with a label
tab and the right-button menu.

This is a large subject and the 1983 sources are the documentation.
\ct{Interface-Framework} and \ct{Interface-Support} in the Browser are where
to read.
